\section{Resultados}

Resultados: Registo dos dados obtidos em tabela (exceto medidas isoladas).

\subsection{Segunda montagem: Lei das malhas}

Procedeu-se à segunda montagem, através da qual se procurou verificar a lei das malhas. A montagem é constituída pelas mesmas resistências da montagem anterior, $R_1 = 1$ k$\Omega$ e $R_2 = 500$ $\Omega$, associadas em série. A tabela seguinte engloba as medições obtidas, os valores teóricos esperados, assim como o desvio e o erro associados às medições.



\begin{table}[h!]
\centering
\caption{Resultados da Montagem 1.}
\begin{tabular}{|c|c|c|c|c|}
\hline
\textbf{Medição} &
\textbf{Valor experimental} &
\textbf{Valor teórico} &
\textbf{Desvio} &
\textbf{Erro (\%)} \\
\hline
$I$ & $29,4 \pm 0.01$ mA & $10,00$ mA &$+19,4$ mA & $194?$ \\
\hline
$V_0$ & $15,18 \pm 0.01$ V & 15,00 V & $+0,18$ V &$1,20$ \\
\hline
$V_1$ & $10,10 \pm 0.01$ V & 10,00 V & $ +0,10$ V &$1,00$ \\
\hline
$V_2$ & $5,06 \pm 0.01$ V & 5,00 V & $ +0,06$ V &$1,20$ \\
\hline
\end{tabular}
\end{table}

\subsubsection{Verificação:}
De acordo com a secção (1.1), temos que:
\begin{equation} 
\Sigma V = 0
\end{equation}

Recorrendo aos valores da Tabela 1:
\begin{equation*} 
(15,18\pm0,01) - (10,10\pm0,01) - (5,06\pm0,01) = (0,02\pm0,03)
\end{equation*}

Visto que o valor experimental $(0,00)$ V está inserido na gama de valores determinados pelas medições e pela incerteza experimental associada, conclui-se que a lei das malhas de $Kirchhoff$ foi verificada.


